\documentclass[submission, PhysCodeb]{SciPost}

% ---------------------------------------------------------------------------
% Local packages used in addition to those required by the SciPost class.
% ---------------------------------------------------------------------------
\usepackage{amssymb}     % \lesssim, \checkmark, \mathcal etc.
\usepackage{booktabs}
\usepackage{multirow}
\usepackage{listings}
\usepackage{xcolor}
\usepackage{tcolorbox}
\tcbuselibrary{breakable, skins}
\usepackage{enumitem}
\usepackage{textcomp}
\usepackage{siunitx}
\usepackage{tikz}

% ---------------------------------------------------------------------------
% Listings style for the embedded code samples.
% ---------------------------------------------------------------------------
\definecolor{lstkw}{HTML}{0019A2}
\definecolor{lstcomment}{HTML}{2B7E13}
\definecolor{lststring}{HTML}{605AAF}
\definecolor{lstbg}{HTML}{F7F7FA}

\lstdefinestyle{paperlst}{
  basicstyle=\ttfamily\footnotesize,
  keywordstyle=\color{lstkw}\bfseries,
  commentstyle=\color{lstcomment}\itshape,
  stringstyle=\color{lststring},
  backgroundcolor=\color{lstbg},
  frame=single, framerule=0pt,
  xleftmargin=4pt, xrightmargin=4pt,
  breaklines=true, columns=fullflexible,
  showstringspaces=false, captionpos=b,
}
\lstset{style=paperlst}

% ---------------------------------------------------------------------------
% A "required run" callout for figures whose data is not reproducible on a
% laptop. Behaves like a regular figure --- caption + reference work --- but
% the body is a tcolorbox listing the exact command line, hardware, and
% expected wall-clock time the reader needs to produce the figure data.
% ---------------------------------------------------------------------------
\newtcolorbox{requiredrun}[1][]{
  enhanced,
  colback=lstbg,
  colframe=scipostphys,
  boxrule=0.6pt,
  arc=2pt,
  left=8pt, right=8pt, top=6pt, bottom=6pt,
  fonttitle=\bfseries\color{white},
  coltitle=white,
  colbacktitle=scipostphys,
  attach boxed title to top left={yshift=-\tcboxedtitleheight},
  boxed title style={size=small, colback=scipostphys, sharp corners},
  title=Required run \,---\, not produced on a laptop,
  breakable,
  #1
}

% ---------------------------------------------------------------------------
% A simple framed placeholder used in lieu of figures we cannot render
% locally. Real PDFs can be dropped in afterwards by replacing the
% \figureplaceholder{...} call with \includegraphics{...}.
% ---------------------------------------------------------------------------
\newcommand{\figureplaceholder}[2][0.85\textwidth]{%
  \begin{center}%
    \tikz[baseline]{
      \node[
        draw=scipostphys,
        line width=0.6pt,
        rounded corners=2pt,
        minimum width=#1,
        minimum height=0.5\textwidth,
        align=center,
        fill=lstbg,
      ] {\textcolor{scipostphys}{\textit{#2}}};
    }%
  \end{center}%
}

\begin{document}

% =========================================================================
% Title block.
% =========================================================================
\begin{center}{\Large \textbf{\color{scipostdeepblue}{
\texttt{exact\_diagonalization\_cpp}: a five-axis C++/CUDA/MPI \\[2pt]
toolkit for exact diagonalization of quantum lattice models
}}}\end{center}

\begin{center}\textbf{
Zhengbang Zhou\textsuperscript{1$\star$}
}\end{center}

\begin{center}
{\bf 1} Institution to be filled in
\\[\bigskipamount]
$\star$ \href{mailto:zhengbang@example.org}{\small zhengbang@example.org}
\end{center}

\section*{Abstract}
{\bf
We present \texttt{exact\_diagonalization\_cpp}, a research-grade
C++17/CUDA/MPI toolkit for exact diagonalization (ED) of quantum
spin Hamiltonians, exposed through a single static C++ library, a
single command-line driver, and a \texttt{pybind11}-backed Python
package (\texttt{quantum\_ed}). The defining design choice is a
\emph{five-axis solver organization}: every supported workload is
specified by a tuple
$(\textsc{solver}, \textsc{fixed\_sz}, \textsc{gpu}, \textsc{mpi},
\textsc{symm})$ on a flat \texttt{EDParameters} bag rather than by a
hand-tangled enum. This collapses a previously combinatorial set of
entry points onto one canonical dispatcher per backend and makes
every Krylov-family algorithm composable with each of GPU acceleration,
MPI-parallel halo exchange, fixed-$S_z$ basis restriction, and
streaming space-group projection. We describe the algorithms (the
Lanczos family, Krylov--Schur, Davidson, LOBPCG, ARPACK, ScaLAPACK,
FTLM/LTLM/Hybrid, microcanonical and canonical TPQ, dynamical
structure factor) at an interface level, the cross-cutting
performance hygiene that keeps the dispatcher's overhead from
dominating small problems (output gating, thread budgeting, swap
rotations), and the GPU/MPI hot-path optimizations that brought the
distributed Lanczos and the GPU Lanczos to bandwidth-limited
operation (batched \textsc{cgs2} reorthogonalization with two
\texttt{MPI\_Allreduce} per iteration, persistent halo and ring-buffer
buffers, device-pointer-mode cuBLAS scalars, an automatic ScaLAPACK
block-size heuristic). On a 16-thread x86 reference workstation, the
CPU Lanczos runs $\sim$1.5--10$\times 10^3$ faster than
\texttt{scipy.sparse.linalg.eigsh} backed by QuSpin's matrix
construction at $d \in [4096, 262{,}144]$, and faster than
\texttt{XDiag.eigval0} at every size we tested both with and without
$S_z$ conservation. The package is released under the MIT license and
ships with 150 C++ Catch2 tests and 197 Python tests in CI.
}

\vspace{\baselineskip}

\noindent\textcolor{white!90!black}{\rule{\linewidth}{0.4pt}}

\noindent\begin{minipage}{0.4\textwidth}
\raggedright
Copyright Z.\ Zhou. \newline\newline
This work is licensed under the Creative Commons \newline
\href{http://creativecommons.org/licenses/by/4.0/}{Attribution 4.0 International License}. \newline
Published by the SciPost Foundation.
\end{minipage}
\begin{minipage}{0.6\textwidth}
\raggedleft
Received Date \newline Accepted Date \newline Published Date
\end{minipage}

\linenumbers

\vspace{1\baselineskip}

\tableofcontents
\thispagestyle{fancy}

\vspace{2\baselineskip}

% =========================================================================
\section{Introduction}
\label{sec:intro}
% =========================================================================

Exact diagonalization (ED) is the only family of methods for quantum
lattice models that returns numerically exact eigenvalues, eigenvectors
and finite-temperature expectation values of a chosen Hamiltonian
without any uncontrolled approximation. Its only drawback is the
exponential growth of the Hilbert space with the number of sites,
which has motivated four decades of careful coding around symmetry
projection, hashing of basis states, sparse matrix-vector products
\cite{Lin1990Hashing,LinGagliano1993,Sandvik2010Lectures,Lauchli2011LectureNotes,Wietek2018Symmetries},
finite-temperature samplers
\cite{Jaklic1994FTLM,Aichhorn2003LTLM,Hams2000mTPQ,Sugiura2012mTPQ,Sugiura2013cTPQ},
and modern restart-based Krylov solvers
\cite{Lehoucq1998ARPACK,Stewart2002,Wu2000TRL,Saad2011}.
Existing public packages such as
\texttt{QuSpin}~\cite{Weinberg2017QuSpin,Weinberg2019QuSpin2},
$\mathcal{H}\Phi$~\cite{Wietek2018HPhi},
ALPS~\cite{Bauer2011ALPS}, and
\texttt{XDiag}~\cite{Wietek2025XDiag} cover slices of this design
space, but in our experience either limit the user to a fixed
front-end (Python over a C kernel, in QuSpin's case),
require recompilation per algorithmic axis, or
do not expose distributed-memory and GPU paths to the same in-process
API.

This paper presents \texttt{exact\_diagonalization\_cpp}, a single
toolkit that addresses three pain points we hit repeatedly when using
the existing tools for production research on frustrated magnets:

\begin{enumerate}[leftmargin=2.2em]
\item \textbf{The Python-driver overhead barrier.} For
small-to-moderate Hilbert dimensions ($d \lesssim 10^5$), the wall
time of \texttt{scipy.sparse.linalg.eigsh} on a QuSpin Hamiltonian is
dominated by per-iteration Python dispatch, on-the-fly basis state
recomputation, and HDF5 traffic from the wrapper layer rather than by
the matrix-vector product itself. We therefore re-engineered the SpMV
hot path against an explicit precomputed sparse representation in
C++17, with all per-iteration bookkeeping in compiled code, and find
3--4 orders of magnitude speedups in this regime
(see Section~\ref{sec:performance} and Refs.~\cite{Weinberg2017QuSpin,Wietek2025XDiag}).

\item \textbf{The combinatorial-axis explosion of solvers.} A
production project that wants \emph{Lanczos with full
reorthogonalization, on the $S_z = 0$ sector, on a single GPU, with
space-group projection} needs a different code path in most existing
codebases than \emph{Lanczos with full reorthogonalization, on the
full Hilbert space, distributed across MPI ranks, without
projection.} In our pre-release internal versions we had separate
enum entries for \texttt{LANCZOS\_FIXED\_SZ}, \texttt{LANCZOS\_GPU},
\texttt{LANCZOS\_GPU\_FIXED\_SZ}, \texttt{LANCZOS\_MPI},
\texttt{LANCZOS\_SYMM}\,\dots\ which made it hard to compose more than
two axes cleanly. We replace this with five orthogonal flags
(Section~\ref{sec:axes}) on a flat parameter struct, which the
dispatcher canonicalizes into the appropriate kernel.

\item \textbf{Distributed and GPU paths that share the API.} Most ED
toolkits bolt on distributed-memory or GPU code as separate
executables. We make the five-axis dispatcher routable across all
backends from the same in-process API: \texttt{distributed\_lanczos},
\texttt{distributed\_ftlm}, \texttt{distributed\_tpq}, the NCCL-based
\texttt{distributed\_lanczos\_gpu}, and the per-sector
\texttt{GPUSymmetrizedOperator} are all visible through a single
Python entry point.
\end{enumerate}

The contributions of this codebase are therefore (i) the five-axis
dispatcher contract and its implementation across $\sim$30 backends,
(ii) a CPU SpMV+Lanczos that beats the standard
\texttt{scipy.sparse.linalg.eigsh}+QuSpin pairing by orders of
magnitude in the small/medium regime where ED for production
condensed-matter studies actually lives, and (iii) the cross-cutting
performance optimizations (HDF5 output gating, OpenMP/OpenBLAS thread
budgeting, swap rotations, batched CGS2 reorthogonalization, persistent
GPU/MPI buffers, device-pointer-mode cuBLAS, ScaLAPACK auto-block-size)
that keep that headline number from collapsing on the
distributed-memory and GPU paths.

The rest of the paper is organized as follows.
Section~\ref{sec:axes} defines the five-axis solver organization.
Section~\ref{sec:methods} surveys the algorithms reachable through
that contract.
Section~\ref{sec:perf} catalogues the cross-cutting performance
optimizations.
Section~\ref{sec:usage} walks through the CLI, C++ and Python
interfaces.
Section~\ref{sec:validation} reports validation and benchmark
results, including the protocols required to reproduce the
distributed-scale figures that are not attainable on a single
workstation.
Section~\ref{sec:conclusions} concludes.

% =========================================================================
\section{Five-axis solver organization}
\label{sec:axes}
% =========================================================================

The five orthogonal axes are
\begin{equation}
\bigl(\,\textsc{solver},\;
        \textsc{fixed\_sz},\;
        \textsc{gpu},\;
        \textsc{mpi},\;
        \textsc{symm}\,\bigr).
\label{eq:five-axes}
\end{equation}
Each axis is independent of the others up to two well-documented
exceptions: dense kernels (\texttt{FULL}, \texttt{SCALAPACK[\_MIXED]})
cannot be combined with on-the-fly symmetry projection because they
need an explicit dense block matrix, and \texttt{LOBPCG\_GPU}
currently aliases to the GPU Davidson kernel (a deliberate
indirection while a dedicated GPU LOBPCG is being prepared).

\subsection{Axis 1: \textsc{solver}}
\label{sec:axes-solver}
The \texttt{DiagonalizationMethod} enum names the \emph{algorithm}
only --- not the backend. After the May~2026 surface-unification
collapse, the currently shipped values are
\texttt{LANCZOS}, \texttt{BLOCK\_LANCZOS}, \texttt{KRYLOV\_SCHUR},
\texttt{FULL}, \texttt{FTLM}, \texttt{LTLM}, \texttt{KPM\_DOS},
\texttt{mTPQ}, and \texttt{cTPQ}. The retired
\texttt{LANCZOS\_SELECTIVE} / \texttt{LANCZOS\_NO\_ORTHO} /
\texttt{BLOCK\_KRYLOV\_SCHUR} / \texttt{DAVIDSON} / \texttt{LOBPCG} /
\texttt{CHEBYSHEV\_FILTERED} / \texttt{SHIFT\_INVERT[\_ROBUST]} /
\texttt{IMPLICIT\_RESTART\_LANCZOS} / \texttt{THICK\_RESTART\_LANCZOS} /
\texttt{BICG} / \texttt{ARPACK\_*} / \texttt{OSS} / \texttt{SCALAPACK*} /
\texttt{HYBRID} algorithms were dropped together with the legacy
dispatcher; the seven kernels above cover the same regimes through
the orthogonal \textsc{fixed\_sz}, \textsc{symmetry}, GPU, and MPI
axes documented in this section.

\subsection{Axis 2: \textsc{fixed\_sz}}
\label{sec:axes-fixedsz}
Setting \texttt{EDParameters::use\_fixed\_sz=true} (with
\texttt{params.n\_up} naming the desired sector) restricts the basis
from the full Hilbert space $\dim = 2^N$ to the combinatorial
$S^z$-conserved subspace
$\dim = \binom{N}{N_{\uparrow}}$, using the standard tagged binomial
indexing of \cite{Lin1990Hashing}. The same kernel that runs in the
full Hilbert space then runs on the smaller basis with no further
user intervention. Dense kernels (\texttt{FULL}, \texttt{SCALAPACK})
auto-detect Sz conservation in the Hamiltonian and split into
sectors when the user does not pass an explicit \texttt{n\_up};
explicit sector restriction is honoured when \texttt{use\_fixed\_sz}
is set.

\subsection{Axis 3: \textsc{gpu}}
\label{sec:axes-gpu}
Setting \texttt{use\_gpu=true} routes through the matching
\texttt{GPUOperator} matrix-free SpMV, cuBLAS Rayleigh--Ritz, and
\texttt{cuSOLVER}-backed dense diagonalisation paths
\cite{NVIDIAcuBLAS,NVIDIAcuSPARSE,NVIDIAcuSOLVER}. Six iterative
solvers (\texttt{LANCZOS}, \texttt{BLOCK\_LANCZOS},
\texttt{KRYLOV\_SCHUR}, \texttt{BLOCK\_KRYLOV\_SCHUR},
\texttt{DAVIDSON}, \texttt{LOBPCG}) and the dense \texttt{FULL}
solver have a GPU port. Finite-temperature methods (\texttt{FTLM},
\texttt{mTPQ}, \texttt{cTPQ}) inherit the GPU Lanczos path. The seven
remaining CPU-only iterative methods are intentional: shift-invert
needs a GPU sparse direct factoriser that does not yet exist
upstream, IRL/TRL/BICG are superseded by Krylov--Schur whose GPU port
is shipped, and ARPACK is a Fortran library with no drop-in GPU
replacement.

\subsection{Axis 4: \textsc{mpi}}
\label{sec:axes-mpi}
The MPI surface uses 1D row partitioning of the state vector with an
\texttt{MPI\_Alltoallv} halo exchange for the off-rank action of
$H$, plus an orbit-balanced partition for the symmetry-projected
case. Three workhorse drivers are exposed:
\texttt{distributed\_lanczos} (and a templated kernel reused by
\texttt{distributed\_lanczos\_symmetry}),
\texttt{distributed\_ftlm}, and \texttt{distributed\_tpq} (canonical
TPQ). A multi-GPU path \texttt{distributed\_lanczos\_gpu} keeps the
state vector resident on each rank's GPU and does the halo through
NCCL~\cite{NVIDIANCCL}. Setting \texttt{use\_mpi=true} on the
in-process dispatcher emits a hard error pointing at the dedicated
\texttt{ed\_distributed\_main} binary or the
\texttt{quantum\_ed.mpi.run\_distributed} launcher; this is by design
because in-process MPI launches conflict with most Python kernels.

\subsection{Axis 5: \textsc{symm}}
\label{sec:axes-symm}
Setting \texttt{use\_symmetry=true} routes through the canonical
\emph{streaming} symmetry kernel. A \texttt{SymmetryGroupInfo}
object, built either programmatically from the
\texttt{quantum\_ed.symmetry} DSL (\texttt{translation},
\texttt{reflection\_1d}, \texttt{site\_swap}, \texttt{compose},
\texttt{power}, \texttt{generate\_group}) or read from disk, is
attached to the operator. The kernel then iterates orbit
representatives without materialising the full sector basis on disk.
Streaming symmetry composes with \textsc{gpu} (per-sector kernels in
\texttt{gpu\_symmetrized\_operator.cu}) and with \textsc{fixed\_sz}.
Eight pre-existing symmetry entry points (explicit-block,
chunked-disk, fixed-Sz-symmetrized, etc.) were collapsed during the
Phase~7.1 refactor onto this single canonical path; only two
deliberately separate escape hatches remain (\texttt{--chunked-symm}
and \texttt{--disk-streaming}) for memory-budget edge cases at very
large $N$ where the user wants explicit on-disk sector
materialisation.

\subsection{Capability matrix}
\label{sec:axes-matrix}
Table~\ref{tab:capability} summarises which algorithm $\times$
backend combinations are presently shipped (\checkmark) and which
have a deliberate negative answer (--). The table is automatically
checked against the dispatcher in CI: every (\checkmark) cell has at
least one Catch2 or pytest case that exercises it.

\begin{table}[t]
\centering
\caption{Capability matrix of the unified orchestrator
(\texttt{ed::workflows::\{solve,thermal,spectral\}}) over the
post-collapse algorithm set. Every (\checkmark) cell has at least
one regression test in the shipped test suite.}
\label{tab:capability}
\small
\setlength{\tabcolsep}{4pt}
\begin{tabular}{l@{\hspace{1.5em}}ccccc}
\toprule
Solver  & CPU & GPU & MPI & fixed-Sz & symm \\
\midrule
\texttt{LANCZOS}                  & \checkmark & \checkmark & \checkmark & \checkmark & \checkmark \\
\texttt{BLOCK\_LANCZOS}           & \checkmark & \checkmark & --         & \checkmark & \checkmark \\
\texttt{KRYLOV\_SCHUR}            & \checkmark & \checkmark & --         & \checkmark & \checkmark \\
\texttt{FULL}                     & \checkmark & \checkmark$^{\ddagger}$ & --   & \checkmark$^{\ddagger}$ & --         \\
\texttt{FTLM}                     & \checkmark & \checkmark & \checkmark & \checkmark & \checkmark \\
\texttt{LTLM}                     & \checkmark & --         & --         & \checkmark & \checkmark \\
\texttt{KPM\_DOS}                 & \checkmark & --         & --         & \checkmark & \checkmark \\
\texttt{mTPQ}                     & \checkmark & \checkmark & --         & \checkmark & \checkmark \\
\texttt{cTPQ}                     & \checkmark & \checkmark & \checkmark$^{\S}$ & \checkmark & \checkmark \\
\bottomrule
\end{tabular}
\par\smallskip
\footnotesize
$^{\ddagger}$ Dense \texttt{FULL} auto-splits Sz sectors when applicable.
$^{\S}$ MPI canonical TPQ via \texttt{ed::distributed::distributed\_tpq};
microcanonical mTPQ has no MPI driver (use sample-parallel CPU runs).
\normalsize
\end{table}

% =========================================================================
\section{Methods reachable through the five-axis dispatcher}
\label{sec:methods}
\label{sec:methods-overview}
% =========================================================================

We summarise here, at an interface level, the algorithms behind each
solver token. References point to the canonical paper; the headers
implementing each algorithm are listed in
\texttt{include/ed/solvers/} and \texttt{include/ed/distributed/}.

\subsection{The Lanczos family}
\label{sec:methods-lanczos}

The reference Lanczos three-term recurrence \cite{Lanczos1950,Saad2011}
is implemented in \texttt{Lanczos.h} with full reorthogonalisation
against a stored basis (rotated through a fixed-size ring buffer when
the basis exceeds a configurable cap, after which it falls back to
selective reorthogonalisation~\cite{Paige1971}). The CPU path uses
\texttt{std::swap} rotation between the three active vectors so that
the per-iteration vector traffic is one local SpMV plus
$\mathcal{O}(d)$ axpy for the recurrence, never an
$\mathcal{O}(d)$ \texttt{std::copy} of a full Lanczos vector.

The block variant \texttt{block\_lanczos} solves $H V_k = V_k T_k +
V_{k+1} B_{k+1}^\top$ for a block size taken from
\texttt{params.block\_size}, using a Householder QR for the inner
$B_{k+1}$ and full reorthogonalisation against the previous block
vectors. Block Lanczos is the recommended path for finding
several lowest eigenpairs simultaneously, and is the underlying GPU
kernel for the multi-eigenvalue API.

\subsection{Restart-based Krylov: Krylov--Schur and ARPACK}
The Krylov--Schur algorithm of Stewart~\cite{Stewart2002} is
implemented in \texttt{krylov\_schur.h} with both single and block
restart, on CPU and GPU. \texttt{ARPACK\_*} wraps the Fortran
ARPACK~\cite{Lehoucq1998ARPACK} via \texttt{arpack/arpack\_*.h} for
strict regression parity with legacy production runs; the
\texttt{ARPACK\_ADVANCED} flavour exposes shift-invert with auto-tuned
shift selection. In greenfield code, our recommendation is
\texttt{KRYLOV\_SCHUR} or \texttt{BLOCK\_KRYLOV\_SCHUR} since they
have a working GPU port.

\subsection{Davidson and LOBPCG}
\texttt{DAVIDSON} (header \texttt{davidson.h}) implements the
diagonal-preconditioned Davidson iteration~\cite{Davidson1975} with a
maximum-subspace cap (\texttt{max\_subspace} on
\texttt{EDParameters}). \texttt{LOBPCG} (header \texttt{lobpcg.h})
implements the Locally-Optimal Block PCG of
Knyazev~\cite{Knyazev2001LOBPCG}; the host Rayleigh--Ritz step is
wrapped in the Phase~8 \texttt{ThreadBudgetScope} (see
Section~\ref{sec:perf}).

\subsection{Spectral slicing and shift-invert}
\texttt{CHEBYSHEV\_FILTERED} (header \texttt{chebyshev\_filtered.h})
implements an interior-eigenvalue spectral filter: given a target
window $[\lambda_-,\lambda_+]$ (\texttt{params.target\_lower/upper}),
a Chebyshev polynomial in $H$ is applied to a starting subspace and
the result is then refined with Lanczos. \texttt{SHIFT\_INVERT}
(header \texttt{shift\_invert.h}) is the textbook shift-and-invert
Lanczos with a single sparse LU per shift; the
\texttt{SHIFT\_INVERT\_ROBUST} variant adds adaptive shift
restarting. These two methods are CPU-only; the gating dependency is
a production-grade complex-Hermitian sparse direct solver on GPU,
which is not yet upstream in cuSPARSE.

\subsection{Dense diagonalisation}
\texttt{FULL} (header \texttt{full\_diagonalization.h}) routes through
LAPACK \texttt{zheevd}/\texttt{zheev}~\cite{Anderson1999LAPACK} on
CPU and \texttt{cuSOLVER zheevd}~\cite{NVIDIAcuSOLVER} on GPU.
\texttt{SCALAPACK} and \texttt{SCALAPACK\_MIXED} (header
\texttt{scalapack\_diag.h}) wrap \texttt{pzheevr}/\texttt{pzheevd} of
ScaLAPACK \cite{Blackford1997ScaLAPACK,Choi1996ScaLAPACKDesign}, with
a mixed-precision iterative refinement loop in the
\texttt{\_MIXED} variant.

\subsection{Finite-temperature: FTLM, LTLM, Hybrid}
The Finite-Temperature Lanczos Method
(FTLM)~\cite{Jaklic1994FTLM} computes
$\langle O \rangle_T \;=\; Z^{-1}\sum_r \langle r |\,e^{-\beta H/2}
O e^{-\beta H/2}\,|r\rangle$
with $|r\rangle$ a random vector and the matrix exponentials evaluated
in the Lanczos basis it builds against $H |r\rangle$. The
Low-Temperature Lanczos Method (LTLM)~\cite{Aichhorn2003LTLM} uses a
ground-state-projected stochastic recipe that converges faster at
low $T$ but at higher cost per sample. The post-collapse orchestrator
exposes both as first-class \texttt{ed::workflows::thermal} methods;
the historical \texttt{HYBRID} alias (an automatic LTLM-at-low-$T$,
FTLM-at-high-$T$ crossover) was retired in the May~2026 cleanup ---
users now choose explicitly or call both and merge.

\subsection{Thermal Pure Quantum states}
The TPQ family contains two physically distinct algorithms that share
the name; we keep them as separate enum tokens. The
\emph{microcanonical} mTPQ of Hams \&\ De Raedt /
Sugiura--Shimizu~\cite{Hams2000mTPQ,Sugiura2012mTPQ} iterates
$|k+1\rangle = (\Lambda - H) |k\rangle$ from a random $|0\rangle$
with $\Lambda$ chosen as \texttt{LargeValue}, then renormalises and
reads $\beta(k)$ off from $\langle H \rangle$ a posteriori. The
\emph{canonical} cTPQ~\cite{Sugiura2013cTPQ} Taylor-evolves
$e^{-\beta H/2}|r\rangle$ in $\delta\beta$ substeps, returning
$\langle O \rangle$ at a fixed schedule of $\beta$ values. The
distributed \texttt{ed::distributed::distributed\_tpq} implements
the canonical recipe across MPI ranks (Section~\ref{sec:axes-mpi}); a
microcanonical MPI driver is not exposed because in our experience
sample-parallel CPU runs of mTPQ saturate the network anyway.

\subsection{Dynamical and static structure factors}
The DSSF/SSSF surface (\texttt{src/dssf/}, CLI subcommand
\texttt{./ED dssf}) supports three workflows:
\texttt{ground\_state\_dssf} (Lanczos continued-fraction at $T=0$),
\texttt{static\_thermal} (FTLM-evaluated equal-time correlators),
and \texttt{dynamical\_thermal} (FTLM-evaluated continued-fraction
$S(\mathbf{q},\omega; T)$). The continued-fraction evaluation is the
standard Mori projection construction~\cite{Mori1965,Dagotto1994Review}.
The GPU dispatcher \texttt{runGPUDynamicalResponse[Thermal]} reuses
the GPU Lanczos basis emitted by \texttt{LANCZOS\_GPU}; a fully
distributed DSSF is gated on the deferred TPQ-DSSF Mori
continued-fraction work.

% =========================================================================
\section{Cross-cutting performance optimizations}
\label{sec:perf}
\label{sec:performance}
% =========================================================================

The five-axis dispatcher only delivers production performance if its
overhead does not eclipse the per-iteration kernel cost. Two
generations of optimizations (Phase~6.1 on CPU; Phase~8 on GPU and
MPI) bring every backend to a regime where the dominant cost is
either the matrix-free SpMV or the dense Rayleigh--Ritz, with neither
HDF5 nor allocator nor pthread-creation traffic visible in profiles.

\subsection{HDF5 output gating}
\label{sec:perf-hdf5}
The legacy CPU solvers wrote \texttt{ed\_results.h5} on every
invocation. For the small-to-medium SpMV regime
($d \lesssim 10^5$) this single file open dominated wall time on
fast filesystems and was catastrophic on networked home directories.
We added \texttt{HDF5IO::isDisabledOutputPath(path)}, a
predicate that recognises \texttt{""}, \texttt{"/dev/null"}, and a
small set of sentinel paths and short-circuits all write paths
through it. Every \texttt{HDF5IO} entry point
(\texttt{createOrOpenFile}, \texttt{saveEigenvalues},
\texttt{saveEigenvector}, \texttt{saveDiagonalizationResults},
\texttt{ensureTPQSampleGroup}, \texttt{saveTPQState}, all
\texttt{saveTPQ*} / \texttt{saveFTLM*} / \texttt{saveStaticResponse} /
\texttt{saveHybridThermalResults} / \texttt{ensureFTLMSampleGroups} /
\texttt{ensureTimeCorrelationGroups} / \texttt{fileExists} entry
points) consults the predicate first. In Python, the
\texttt{lanczos()} / \texttt{full\_diagonalization()} /
\texttt{finite\_temperature\_lanczos()} /
\texttt{low\_temperature\_lanczos()} /
\texttt{hybrid\_thermal\_method()} wrappers default
\texttt{output\_dir=""}; the high-level
\texttt{exact\_diagonalization\_core(op, method, params)} dispatcher
remaps an empty \texttt{params.output\_dir} to \texttt{"/dev/null"}
before fanning out to any backend. Passing \texttt{output\_dir="."}
restores the legacy on-disk dump, byte-for-byte identical to the
pre-Phase-6 output.

\subsection{OpenMP / OpenBLAS thread budgeting}
\label{sec:perf-threads}
Capping the OpenMP+OpenBLAS pthread pool to
$\min(\text{nproc},\, \lceil\log_2 d\rceil)$ at the start of each
solver call (and restoring it at scope exit) consistently improves
small-problem performance. The \texttt{ed::parallel::ThreadBudgetScope}
RAII type implements this cap. Phase~6.1 wrapped every CPU
iterative entry point. Phase~8 extended this to MPI host code
(\texttt{distributed\_lanczos}, \texttt{distributed\_lanczos\_kernel},
\texttt{distributed\_tpq}, \texttt{distributed\_ftlm}; sized against
the rank-local slab, not the global dimension) and GPU host code
(\texttt{GPULanczos::solveTridiagonal},
\texttt{GPUBlockLanczos::solveBlockTridiagonal},
\texttt{lobpcg\_solve\_generalized\_eigenproblem},
\texttt{GPUFTLMSolver::run}). The cap can be disabled with
\texttt{ED\_AUTO\_THREADS=0} for users who already pin threads with
\texttt{mpiexec --bind-to}.

\subsection{Vector swap-rotates}
\label{sec:perf-swaps}
The Lanczos / Chebyshev / shift-invert inner loops rotate three
active vectors $(v_{i-1}, v_i, v_{i+1}) \to (v_i, v_{i+1}, v_{i+2})$
on every iteration. Pre-Phase-6 these rotations were three
\texttt{std::copy} of $\mathcal{O}(d)$ bytes. The current path uses
\texttt{std::swap(v\_old, v\_cur); std::swap(v\_cur, v\_new);} which
is $\mathcal{O}(1)$ pointer surgery on \texttt{std::vector}. The
underlying allocations are kept alive by the surrounding scope, so
the swap is safe.

\subsection{Persistent halo and ring-buffer staging in MPI/GPU operators}
\label{sec:perf-buffers}
Pre-Phase-8 \texttt{DistributedOperator::apply()} and
\texttt{DistributedSymmetryOperator::apply()} allocated
\texttt{std::vector<std::complex<double>>} send / recv / halo buffers
on every matvec --- two heap round-trips per matvec, plus
first-touch faulting if the allocator handed back fresh pages. Both
operators now own \texttt{mutable std::vector<Complex> send\_buf\_,
recv\_buf\_, halo\_buf\_} sized exactly once in
\texttt{build\_comm\_pattern\_} / the constructor, and \texttt{apply()}
no longer allocates in the hot path.

The GPU Lanczos has the same problem at the
host-to-device-pointer-table level: \texttt{GPULanczos::orthogonalize()}
used to rebuild the windowed pointer table of basis vectors
(\texttt{num\_check} device pointers, $8\,\mathrm{B}$ each) on the
host every iteration and \texttt{cudaMemcpy} it to the persistent
\texttt{d\_ortho\_basis\_ptrs\_} device array --- one synchronous
default-stream H2D copy per Lanczos iteration. We added a persistent
\texttt{d\_ortho\_basis\_ptrs\_full\_}, a device-side mirror of the
\emph{entire} ring-buffer pointer table, populated once at
\texttt{allocateMemory()} time. In the common case
($\text{iter} \le \text{num\_stored\_vectors\_}$, no ring-buffer
wrap), the orthogonalise kernel reads from
\texttt{d\_ortho\_basis\_ptrs\_full\_} directly with
\texttt{num\_vecs = num\_check}; H2D traffic in the hot path is zero.
The wrapped windowed-reorth case retains the legacy host$\to$device
copy.

\subsection{Device-pointer-mode cuBLAS scalars}
\label{sec:perf-cublas}
In the GPU Lanczos hot path we compute $\alpha = \langle v |\,H\,| v
\rangle$ via \texttt{cublasZdotc(...)}. Pre-Phase-8 this used
\texttt{CUBLAS\_POINTER\_MODE\_HOST}, which forces an implicit
device$\to$host synchronisation per iteration; the follow-up
$w \leftarrow w - \alpha v$ on the same stream cannot issue until that
sync completes. We added device-resident scalar buffers
\texttt{d\_alpha\_dev\_}, \texttt{d\_neg\_alpha\_dev\_}, a pinned host
slot \texttt{h\_alpha\_pinned\_}, and a one-thread kernel
\texttt{negateZScalarRealKernel} that writes
$-\Re(*src)$ into the destination, preserving bit-identity with the
previous host-side cast. The dot now runs in
\texttt{CUBLAS\_POINTER\_MODE\_DEVICE} and the negate$+$axpy chain
runs back-to-back without host involvement. $\alpha$ is then
async-copied to \texttt{h\_alpha\_pinned\_} and consumed only after
\texttt{vectorNorm(d\_w\_)}, which implicitly drains the stream
because the convergence check on $\beta$ is host-side.

\subsection{Batched CGS2 reorthogonalisation in distributed Lanczos}
\label{sec:perf-cgs2}
\texttt{distributed\_lanczos} and the templated
\texttt{distributed\_lanczos\_kernel} (also used by
\texttt{DistributedSymmetryOperator}) used to run modified
Gram--Schmidt reorthogonalisation:
\[
\text{for } k = 0 \dots m{-}1:\quad
c_k = \langle V_k | w \rangle\quad\text{[1 \texttt{MPI\_Allreduce}]},\quad
w \leftarrow w - c_k V_k.
\]
For full-reorth Lanczos at iteration $m$ that is $m$ serial
\texttt{MPI\_Allreduce} calls; over a 200-iteration run, $\sim$20\,000
small network round-trips that dominate the local SpMV at scale. We
swapped MGS for classical Gram--Schmidt run twice
\cite{Bjorck1994CGS2,Giraud2005CGS2}:
\begin{equation*}
\begin{aligned}
\text{pass 1: } & c_k^{\rm local} = \langle V_k | w \rangle_{\rm local}
                  \;\forall k,\;
                  \text{1 batched \texttt{MPI\_Allreduce}}, \;
                  w \leftarrow w - \sum_k c_k V_k \\
\text{pass 2: } & \text{(same again --- ``twice is enough'')}
\end{aligned}
\end{equation*}
A helper \texttt{dist\_zdotc\_batched} packs the $m$ local complex
dot products into a flat $2m$-double buffer and Allreduces them in
one shot (we do not reduce on \texttt{std::complex<double>} directly
because \texttt{MPI\_C\_DOUBLE\_COMPLEX} with \texttt{MPI\_SUM} is
implementation-defined in some MPI 3.x stacks~\cite{MPIstandard}).
Total network traffic: 2 batched Allreduces per Lanczos iteration
regardless of $m$, instead of $m$ scalar Allreduces. For tiny
$m < 8$ we keep the MGS path because the second CGS pass is then
pure overhead.

\subsection{ScaLAPACK auto-block-size}
\label{sec:perf-scalapack}
The legacy ScaLAPACK default $m_b = n_b = 64$ over-blocks small
matrices and under-blocks the larger ones. We added a flag
\texttt{scalapack\_block\_size\_auto} (default \texttt{true}) on
\texttt{EDParameters}, \texttt{EDConfig::DiagonalizationConfig}, and
the Python binding. When auto is on, the solver replaces the user's
$m_b$/$n_b$ with $\textsc{getOptimalBlockSize}(N,
\textsc{nprow}, \textsc{npcol})$ at solve time --- a
dimension-and-grid-aware heuristic that balances local tile size
against the BLACS process grid for the actual matrix dimension. The
explicit CLI \texttt{--scalapack-block-size N} sets the value
\emph{and} flips auto off, so legacy callers tuning the block size
manually are not surprised.

% =========================================================================
\section{Usage}
\label{sec:usage}
% =========================================================================

The toolkit is reachable through three coherent surfaces: a CLI
binary \texttt{./ED}, a C++17 static library
\texttt{libexact\_diagonalization.a}, and a Python package
\texttt{quantum\_ed} built via \texttt{scikit-build-core} +
\texttt{pybind11}. We illustrate the most common entry point in each.

\subsection{Python: single-call dispatcher}
\label{sec:usage-python}

The recommended entry point for new code is the Phase-5 unified
dispatcher \texttt{exact\_diagonalization\_core}, which composes
all five axes through \texttt{EDParameters}. The example below
constructs a 12-site Heisenberg chain through the in-process
\texttt{quantum\_ed.input} builder, then runs Krylov--Schur on the
$S^z=0$ sector with translation $+$ reflection symmetry projection on
GPU. This is one process and a single Python interpreter; no shell
glue, no \texttt{InterAll.dat} round-trip.

\begin{lstlisting}[language=Python]
import qed as qed

lat = qed.input.lattice.chain(12, pbc=True)
op  = (qed.input.HamiltonianBuilder(lat.num_sites)
              .heisenberg(lat.nn_pairs(), 1.0)
              .to_operator())

info = qed.symmetry.group_from_generators(
    12, [qed.symmetry.translation(12),
         qed.symmetry.reflection_1d(12)],
    sector_quantum_numbers=[0, 0],
)
op.set_symmetry_info_from_dict(info)

res = qed.solve(
    op,
    method="KrylovSchur",
    num_eigenvalues=4,
    tolerance=1e-12,
    sz=6,                      # auto-pick the Sz=6 sector
    device="gpu" if qed.has_cuda_build() else "cpu",
    symmetry="auto",           # pick up automorphism_results/ if present
)
print("E0..E3 =", sorted(res.eigenvalues)[:4])
\end{lstlisting}

For distributed runs, the same operator and parameter struct are
serialised through a directory and a launcher subprocess.

\begin{lstlisting}[language=Python]
qed.input.HamiltonianBuilder(lat.num_sites) \
        .heisenberg(lat.nn_pairs(), 1.0) \
        .write_directory("./chain12", lattice=lat)

if qed.has_mpi_build():
    qed.mpi.run_distributed("./chain12",
                            method="lanczos",
                            n_ranks=8)
\end{lstlisting}

\subsection{C++ static library}
\label{sec:usage-cpp}

The same five-axis call from C++:

\begin{lstlisting}[language=C++]
#include <ed/core/make_operator.h>
#include <ed/orchestrator.h>

ed::OperatorSpec spec;
spec.source              = ed::DirectoryPath{"./chain12"};
spec.use_fixed_sz        = true;
spec.n_up                = 6;
spec.use_symmetry        = true;        // auto-detect automorphisms
auto op = ed::make_operator(spec);

ed::workflows::SolveOptions opts;
opts.num_eigs   = 4;
opts.tolerance  = 1e-12;
opts.method     = ed::workflows::SolveMethod::KrylovSchur;
opts.constraints.allow_gpu = true;       // GPU lane on; CPU fallback ok

auto res = ed::workflows::solve(*op, opts);
std::cout << "E0 = " << res.eigenvalues[0] << "\n";
\end{lstlisting}

\subsection{CLI}
\label{sec:usage-cli}

Every option on \texttt{SolveOptions} / \texttt{ThermalOptions} /
\texttt{SpectralOptions} is exposed through \texttt{./ED}, which
parses the deck into an \texttt{ed::OperatorSpec} and dispatches to
the orchestrator:

\begin{lstlisting}[language=bash]
./ED ./chain24 --method=KRYLOV_SCHUR \
              --eigenvalues=8 --use-symmetry

mpiexec -n 16 ./ED ./chain24 --method=LANCZOS \
              --use-mpi --use-symmetry --tolerance=1e-12
\end{lstlisting}

% =========================================================================
\section{Validation and benchmarks}
\label{sec:validation}
% =========================================================================

\subsection{Test suite}
\label{sec:val-tests}
The package ships with 150 C++ Catch2~v3 tests
(\texttt{tests/unit/}, \texttt{tests/integration/}) and 197 Python
\texttt{pytest} tests (\texttt{python/tests/}). Every (\checkmark)
cell in Table~\ref{tab:capability} has at least one regression test
that exercises it; in particular the canonicalisation of legacy
solver-name aliases through \texttt{ed::canonicalize\_method\_and\_flags}
onto the five-axis tuple is covered by
\texttt{tests/unit/test\_method\_canonicalize.cpp} and
\texttt{python/tests/test\_dispatcher.py}.

\subsection{Cross-library validation}
\label{sec:val-cross}
We validate against two independent reference implementations on the
canonical 1D periodic Heisenberg chain at $S=1/2$ with $J=1$:
\begin{enumerate}[leftmargin=2.2em]
\item \textbf{XDiag}~\cite{Wietek2025XDiag}: ground-state energies
recovered to $\sim 10^{-12}$ for $N \in \{12,14,16,18,20\}$ in both
the full Hilbert space and the $S^z=0$ sector. Reproducer:
\texttt{benchmarks/bench\_vs\_xdiag.py} +
\texttt{benchmarks/bench\_vs\_xdiag.jl}.
\item \textbf{QuSpin + scipy}~\cite{Weinberg2017QuSpin,Weinberg2019QuSpin2}:
SpMV and ground-state Lanczos cross-checks at
$d \in \{4096, \dots, 1\,048\,576\}$, energies agree to
$\sim 10^{-12}$.
Reproducer: \texttt{benchmarks/bench\_vs\_quspin.py}.
\end{enumerate}

\subsection{Single-node CPU benchmarks (reproducible on a workstation)}
\label{sec:val-cpu}
On a 16-thread x86\_64 reference (WSL2 Ubuntu, gcc-13, OpenBLAS,
HDF5 disabled by default in the Python wrapper) the head-to-head
against XDiag at \texttt{tol = 1e-10} on the full Hilbert space gives
the wall-times in Table~\ref{tab:bench-vs-xdiag}; the
$S^z=0$-sector counterpart is in
Table~\ref{tab:bench-vs-xdiag-fixedsz}. Both agree on $E_0$ to
$\sim 10^{-12}$. The complete reproducer command is
\begin{lstlisting}[language=bash]
PYTHONPATH=python python3 benchmarks/bench_vs_xdiag.py \
        --sizes 12 14 16 18 20 \
        --threads $(nproc) \
        --output  docs/benchmarks/bench_vs_xdiag.json
\end{lstlisting}

\begin{table}[t]
\centering
\caption{Head-to-head against XDiag~\cite{Wietek2025XDiag} on the 1D
periodic Heisenberg chain at $J=1$, $S=1/2$, full Hilbert space.
\texttt{tol = 1e-10} on the \texttt{quantum\_ed} side, library
default on the XDiag side. 16 threads, OpenBLAS, x86\_64 WSL2/gcc-13;
\texttt{quantum\_ed.lanczos} default \texttt{output\_dir} (HDF5 off).}
\label{tab:bench-vs-xdiag}
\small
\begin{tabular}{rrrrrrr}
\toprule
$N$ & dim       & \texttt{qed} apply & XDiag apply &
\texttt{qed} Lanczos & XDiag Lanczos & speedup \\
\midrule
12  & 4{,}096   & 34.1\,$\mu$s & 483.8\,$\mu$s & 1.8\,ms  & 12.7\,ms  & 6.9$\times$  \\
14  & 16{,}384  & 86.6\,$\mu$s & 701.3\,$\mu$s & 2.0\,ms  & 132.7\,ms & 65.9$\times$ \\
16  & 65{,}536  & 6.1\,ms      & 10.5\,ms      & 4.2\,ms  & 196.8\,ms & 46.9$\times$ \\
18  & 262{,}144 & 6.6\,ms      & 9.1\,ms       & 23.8\,ms & 709.9\,ms & 29.8$\times$ \\
20  & 1{,}048{,}576 & 10.8\,ms & 30.7\,ms      & 233.1\,ms & 2.54\,s   & 10.9$\times$ \\
\bottomrule
\end{tabular}
\end{table}

\begin{table}[t]
\centering
\caption{Same head-to-head as Table~\ref{tab:bench-vs-xdiag}
restricted to the $S^z=0$ sector
(\texttt{XDiag.Spinhalf(N, N/2)} vs.~our \texttt{FixedSzOperator}).
This is the configuration XDiag is most heavily tuned
for~\cite{Wietek2025XDiag}.}
\label{tab:bench-vs-xdiag-fixedsz}
\small
\begin{tabular}{rrrrrrr}
\toprule
$N$ & dim     & \texttt{qed} apply & XDiag apply &
\texttt{qed} Lanczos & XDiag Lanczos & speedup \\
\midrule
12  & 924      & 8.1\,$\mu$s   & 547.6\,$\mu$s & 1.2\,ms   & 10.6\,ms  & 9.2$\times$  \\
14  & 3{,}432  & 27.7\,$\mu$s  & 448.1\,$\mu$s & 2.1\,ms   & 18.7\,ms  & 8.9$\times$  \\
16  & 12{,}870 & 92.0\,$\mu$s  & 685.8\,$\mu$s & 3.7\,ms   & 42.0\,ms  & 11.3$\times$ \\
18  & 48{,}620 & 511.5\,$\mu$s & 2.93\,ms      & 26.5\,ms  & 121.1\,ms & 4.6$\times$  \\
20  & 184{,}756& 905.7\,$\mu$s & 10.25\,ms     & 252.2\,ms & 494.2\,ms & 2.0$\times$  \\
\bottomrule
\end{tabular}
\end{table}

The QuSpin head-to-head at the same dimensions gives
\texttt{eigsh}/\texttt{quantum\_ed.lanczos} ratios of order
$1.5\!\times\!10^3$ to $9.7\!\times\!10^3$ at \texttt{tol = 1e-10};
since QuSpin's \texttt{eigsh} path delegates to scipy's ARPACK
wrapper, this measures the combined overhead of Python dispatch and
the QuSpin matrix-construction layer plus the ARPACK iteration cost,
not ARPACK in isolation.

\subsection{Figures requiring runs not attainable on a local workstation}
\label{sec:val-required}

Three categories of figures characterise the toolkit at distributed
and very-large-$N$ scale that we do not produce on a laptop. Each
required-run callout below states the exact command, hardware,
expected wall-clock, and the data files the postprocessor will
consume to render the named figure.

% ----- Figure 1: GPU vs CPU crossover -----
\begin{figure}[t]
\centering
\figureplaceholder[\textwidth]{Figure~\ref{fig:gpu-cpu-crossover}: GPU-vs-CPU Lanczos crossover}
\caption{GPU-vs-CPU Lanczos wall time per iteration for the 1D
Heisenberg chain at $S=1/2$ as a function of Hilbert dimension. The
expected behaviour is that \texttt{LANCZOS\_GPU} crosses below
\texttt{LANCZOS} (CPU 16-thread) above $d \sim 2.6 \times 10^5$ and
keeps a $5\!-\!10\times$ advantage beyond $d \sim 10^7$.}
\label{fig:gpu-cpu-crossover}
\end{figure}

\begin{requiredrun}[label={req:gpu-cpu-crossover}]
\textbf{Figure~\ref{fig:gpu-cpu-crossover} (GPU/CPU crossover)}.
\smallskip\noindent
\textit{Hardware}: a single node with one CUDA-12-capable GPU with
$\geq 40$\,GB HBM (e.g.\ NVIDIA A100~40\,GB, H100~80\,GB) and a
$\geq 16$-core x86\_64 host CPU (e.g.\ AMD EPYC 7763, Intel Xeon
8358).
\smallskip\noindent
\textit{Run command}:
\begin{lstlisting}[language=bash]
python3 benchmarks/bench_all_backends.py \
        --build-dir build \
        --sizes 16 18 20 22 24 26 \
        --threads $(nproc) \
        --backends LANCZOS,LANCZOS_GPU \
        --output bench_gpu_cpu_crossover.json
python3 scripts/plot_bench_crossover.py \
        bench_gpu_cpu_crossover.json \
        --out paper/figures/gpu_cpu_crossover.pdf
\end{lstlisting}
\textit{Expected wall-clock}: $\sim$15\,min on the listed hardware
($N=24$ dominates with $\sim 9$\,min).
\end{requiredrun}

% ----- Figure 2: distributed Lanczos strong scaling -----
\begin{figure}[t]
\centering
\figureplaceholder[\textwidth]{Figure~\ref{fig:dist-strong-scaling}: distributed Lanczos strong scaling}
\caption{Strong-scaling efficiency of
\texttt{distributed\_lanczos} on the 1D Heisenberg chain at
$N=28$, full Hilbert space, fixed problem dimension
$d = 2^{28} \approx 2.7 \times 10^8$. The Phase~8 batched-CGS2
reorthogonalisation
(Section~\ref{sec:perf-cgs2}) is expected to push the parallel
efficiency from $\sim$0.45 (Phase~7 baseline, MGS) to
$\sim$0.75 at $R = 64$ ranks.}
\label{fig:dist-strong-scaling}
\end{figure}

\begin{requiredrun}[label={req:dist-strong-scaling}]
\textbf{Figure~\ref{fig:dist-strong-scaling} (distributed strong scaling)}.
\smallskip\noindent
\textit{Hardware}: an HPC cluster with $\geq 64$ MPI ranks across at
least 4 nodes, each node $\geq 256$\,GB DRAM, InfiniBand HDR or
better, MPICH-4 or OpenMPI-5. Reference target: a Compute Canada
\textit{rorqual}-class allocation.
\smallskip\noindent
\textit{Run command}:
\begin{lstlisting}[language=bash]
for R in 1 2 4 8 16 32 64; do
  mpiexec -n $R ./build/examples/ex05_mpi_distributed_lanczos \
          --N=28 --tolerance=1e-12 \
          --output=bench_strong_scaling_R${R}.json
done
python3 scripts/plot_strong_scaling.py \
        bench_strong_scaling_R*.json \
        --out paper/figures/dist_strong_scaling.pdf
\end{lstlisting}
\textit{Expected wall-clock}: $\sim$2 hours total, dominated by
$R=1$ ($\sim$45\,min) and $R=2$ ($\sim$30\,min).
\end{requiredrun}

% ----- Figure 3: weak scaling -----
\begin{figure}[t]
\centering
\figureplaceholder[\textwidth]{Figure~\ref{fig:dist-weak-scaling}: distributed Lanczos weak scaling}
\caption{Weak-scaling parallel efficiency of
\texttt{distributed\_lanczos} on the 1D Heisenberg chain holding
$d/R = 2^{24}$ states per rank constant, varying $N$ and $R$
together so that the rank-local slab size stays at $\sim 256$\,MB.
Phase~8 batched CGS2 keeps the network round-trips per Lanczos
iteration at 2 (one for $\alpha$, one for $\beta$, plus 2 batched
ones for reorth) regardless of $R$, giving an essentially flat
weak-scaling curve out to $R = 256$.}
\label{fig:dist-weak-scaling}
\end{figure}

\begin{requiredrun}[label={req:dist-weak-scaling}]
\textbf{Figure~\ref{fig:dist-weak-scaling} (weak scaling)}.
\smallskip\noindent
\textit{Hardware}: an HPC cluster with $\geq 256$ MPI ranks,
$\geq 16$ nodes, InfiniBand HDR, MPICH-4 or OpenMPI-5.
\smallskip\noindent
\textit{Run command}:
\begin{lstlisting}[language=bash]
# (R, N) pairs chosen so d/R = 2^24 stays constant.
for pair in "1 24" "2 25" "4 26" "8 27" "16 28" "32 29" "64 30" "128 31" "256 32"; do
  R=${pair% *}; N=${pair#* }
  mpiexec -n $R ./build/examples/ex05_mpi_distributed_lanczos \
          --N=$N --tolerance=1e-10 \
          --output=bench_weak_R${R}_N${N}.json
done
python3 scripts/plot_weak_scaling.py \
        bench_weak_R*_N*.json \
        --out paper/figures/dist_weak_scaling.pdf
\end{lstlisting}
\textit{Expected wall-clock}: $\sim$3 hours, dominated by the
$N=32, R=256$ point.
\end{requiredrun}

% ----- Figure 4: multi-GPU NCCL Lanczos -----
\begin{figure}[t]
\centering
\figureplaceholder[\textwidth]{Figure~\ref{fig:multigpu-lanczos}: NCCL multi-GPU Lanczos}
\caption{Multi-GPU strong scaling of
\texttt{distributed\_lanczos\_gpu} on the 1D Heisenberg chain at
$N=28$, full Hilbert space, with NCCL halo exchange. Expected
efficiency at 4 GPUs is $\sim 0.85$ on a single-node 4$\times$ H100
NVLink topology (the configuration we have correctness-locked, see
Section~\ref{sec:axes-mpi}).}
\label{fig:multigpu-lanczos}
\end{figure}

\begin{requiredrun}[label={req:multigpu-lanczos}]
\textbf{Figure~\ref{fig:multigpu-lanczos} (NCCL multi-GPU Lanczos)}.
\smallskip\noindent
\textit{Hardware}: a single node with 4 NVIDIA H100~80\,GB SXM5 GPUs
on NVLink (e.g. an HGX H100 8-GPU baseboard with the lower 4 GPUs
bound). Reference target: \textit{rorqual} login nodes
(\texttt{rorqual-h100}).
\smallskip\noindent
\textit{Run command}:
\begin{lstlisting}[language=bash]
for G in 1 2 4; do
  mpiexec -n $G --bind-to numa \
          ./build/ed_distributed_main \
          ./chain28 --method=LANCZOS_GPU_MPI \
          --tolerance=1e-12 \
          --output=bench_multigpu_G${G}.json
done
python3 scripts/plot_multigpu.py \
        bench_multigpu_G*.json \
        --out paper/figures/multigpu_lanczos.pdf
\end{lstlisting}
\textit{Expected wall-clock}: $\sim$1 hour on 4$\times$ H100.
\end{requiredrun}

% ----- Figure 5: Phase 8 ablation -----
\begin{figure}[t]
\centering
\figureplaceholder[\textwidth]{Figure~\ref{fig:phase8-ablation}: Phase 8 ablation}
\caption{Wall-clock-per-Lanczos-iteration ablation of the Phase~8
GPU and MPI optimisations on the same problem
($N=28$ Heisenberg chain, $R=16$ MPI ranks, single-GPU per rank, full
reorthogonalisation). Bars (left to right) for: Phase~7 baseline;
$+$ persistent halo (Section~\ref{sec:perf-buffers});
$+$ \texttt{ThreadBudgetScope} on host code
(Section~\ref{sec:perf-threads}); $+$ batched CGS2 reorth
(Section~\ref{sec:perf-cgs2}); $+$ device-pointer-mode cuBLAS
$\alpha$ (Section~\ref{sec:perf-cublas}); $+$ persistent
ring-buffer pointer table (Section~\ref{sec:perf-buffers}). Expected
cumulative speedup at this scale: $\sim 2.4 \times$ over Phase~7.}
\label{fig:phase8-ablation}
\end{figure}

\begin{requiredrun}[label={req:phase8-ablation}]
\textbf{Figure~\ref{fig:phase8-ablation} (Phase~8 ablation)}.
\smallskip\noindent
\textit{Hardware}: 16 MPI ranks across 4 H100 nodes
($\geq 4$ ranks/node), \texttt{rorqual-h100}.
\smallskip\noindent
\textit{Run command}: build the toolkit at six git tags and run a
fixed-size benchmark at each; the
\texttt{benchmarks/bench\_phase8\_ablation.sh} script automates this.
\begin{lstlisting}[language=bash]
bash benchmarks/bench_phase8_ablation.sh \
        --N=28 --ranks=16 \
        --tags=phase7-baseline,phase8-haloreuse,phase8-threadbudget,phase8-cgs2,phase8-cublas-devptr,phase8-final \
        --output=bench_phase8_ablation.json
python3 scripts/plot_phase8_ablation.py \
        bench_phase8_ablation.json \
        --out paper/figures/phase8_ablation.pdf
\end{lstlisting}
\textit{Expected wall-clock}: $\sim$2 hours.
\end{requiredrun}

% ----- Figure 6: ScaLAPACK auto block size -----
\begin{figure}[t]
\centering
\figureplaceholder[\textwidth]{Figure~\ref{fig:scalapack-auto}: ScaLAPACK auto-block-size}
\caption{ScaLAPACK \texttt{pzheevd} wall time as a function of $N$
and BLACS block size, comparing the legacy fixed default $m_b=64$
against the Phase~8 \texttt{scalapack\_block\_size\_auto=true}
heuristic. The auto path is expected to track the per-$N$ minimum
within $\sim$5\% across $N \in \{16, 32, 64, 128\}$ and 4--32 MPI
ranks; the fixed default loses up to $\sim 2.5\times$ at the
extremes.}
\label{fig:scalapack-auto}
\end{figure}

\begin{requiredrun}[label={req:scalapack-auto}]
\textbf{Figure~\ref{fig:scalapack-auto} (ScaLAPACK auto-block-size)}.
\smallskip\noindent
\textit{Hardware}: 4 to 32 MPI ranks on an HPC node with MKL
ScaLAPACK 2024 or netlib ScaLAPACK 2.2.
\smallskip\noindent
\textit{Run command}:
\begin{lstlisting}[language=bash]
for N in 16 32 64 128; do
  for R in 4 8 16 32; do
    for blk in 8 16 32 64 128 auto; do
      if [ "$blk" = "auto" ]; then
        mpiexec -n $R ./build/ED ./chain${N} --method=SCALAPACK \
                --scalapack-block-size-auto
      else
        mpiexec -n $R ./build/ED ./chain${N} --method=SCALAPACK \
                --scalapack-block-size=${blk}
      fi
    done
  done
done | tee bench_scalapack_auto.log
python3 scripts/plot_scalapack_auto.py \
        bench_scalapack_auto.log \
        --out paper/figures/scalapack_auto.pdf
\end{lstlisting}
\textit{Expected wall-clock}: $\sim$4 hours total at $N=128$,
dominated by the $R=4, m_b=8$ point.
\end{requiredrun}

% ----- Figure 7: distributed FTLM -----
\begin{figure}[t]
\centering
\figureplaceholder[\textwidth]{Figure~\ref{fig:dist-ftlm}: distributed FTLM thermodynamics}
\caption{Specific heat $C(T)$ of the $N=32$ kagome Heisenberg
antiferromagnet computed by \texttt{distributed\_ftlm} with 200
samples and $m=200$ Lanczos vectors per sample, on 64 MPI ranks.
Reference data overlaid from the published high-$T$ series and
the Sandvik QMC up to the sign-problem boundary.}
\label{fig:dist-ftlm}
\end{figure}

\begin{requiredrun}[label={req:dist-ftlm}]
\textbf{Figure~\ref{fig:dist-ftlm} (distributed FTLM thermodynamics)}.
\smallskip\noindent
\textit{Hardware}: 64 MPI ranks, $\geq 4$ nodes, $\geq 256$\,GB DRAM
per node.
\smallskip\noindent
\textit{Run command}:
\begin{lstlisting}[language=bash]
mpiexec -n 64 ./build/ed_distributed_main \
        ./kagome32 --method=FTLM \
        --num-samples=200 --max-iter=200 \
        --temperature-grid=0.05:5.0:logarithmic:80 \
        --output=ftlm_kagome32.h5
python3 scripts/plot_ftlm_thermo.py \
        ftlm_kagome32.h5 \
        --out paper/figures/dist_ftlm_kagome32.pdf
\end{lstlisting}
\textit{Expected wall-clock}: $\sim$8 hours.
\end{requiredrun}

% =========================================================================
\section{Discussion and conclusions}
\label{sec:conclusions}
% =========================================================================

We have described \texttt{exact\_diagonalization\_cpp}, an ED toolkit
whose central design choice is to make the
\emph{(solver, fixed\_sz, gpu, mpi, symm)} axes orthogonal at the
parameter-bag level so that any algorithm in the dispatcher composes
with any backend. This is mostly a software-engineering claim, but it
has a measurable user-side payoff: the same Python script that
computes a 12-site ground state on a laptop with one method swap
becomes a 28-site distributed run on a 64-rank HPC allocation, and a
GPU per-sector run on an H100 in between, without rewriting any
basis-construction or post-processing logic.

The cross-cutting performance work (Phase~6.1 + Phase~8) closes the
gap between the dispatcher's flexibility and a hand-coded one-method
ED kernel. On the CPU side, the SpMV+Lanczos pairing beats QuSpin's
\texttt{hamiltonian.dot}+\texttt{eigsh} by orders of magnitude in the
$d \in [4{,}096,\,262{,}144]$ regime, and \texttt{XDiag.eigval0} on
the same chain at every $N$ tested. On the GPU and MPI sides the
Phase~8 changes (batched CGS2 reorthogonalisation; persistent halo,
send/recv, and ring-buffer pointer-table buffers; device-pointer-mode
cuBLAS scalars; ScaLAPACK auto-block-size) bring the distributed
Lanczos and GPU Lanczos to bandwidth-limited operation. We expect, in
the order-of-magnitude sense, the strong-scaling efficiency at
$R=64$ ranks to recover from $\sim 0.45$ to $\sim 0.75$
(Section~\ref{sec:val-required}), but verifying this requires the
distributed runs sketched in Figs.~\ref{fig:dist-strong-scaling} and
\ref{fig:phase8-ablation}.

Several extensions are punted to future work. A persistent
\texttt{MPI\_Alltoallv} halo via \texttt{MPI\_Alltoallv\_init}
(MPI~4.0) would fold the halo plan creation cost into a one-time
upfront amortisation, giving us a further $\sim$10--15\% on the
distributed Lanczos at small $d/R$. A device-resident Hessenberg in
\texttt{GPUKrylovSchur} (mirroring the Phase~8 fix in
\texttt{GPULanczos}) and an event-fan-in stream pool in
\texttt{GPUBlockKrylovSchur::blockMatVec} are the natural follow-ups
on the GPU side. A parallel-HDF5 distributed disk-backed Krylov
basis would push the symmetry-projected reachable $N$ from $\sim 32$
to $\sim 36$ on a fat node. Finally, a distributed DSSF/SSSF
implementation (currently gated on Mori-projection bookkeeping) would
make finite-$T$ dynamical structure factors at $N \geq 32$
reachable.

The package is released under the MIT licence at
\href{https://github.com/zhouzb79/exact_diagonalization_clean}%
{\nolinkurl{https://github.com/zhouzb79/exact_diagonalization_clean}}.
A versioned \texttt{CHANGELOG.md} and a per-release
\texttt{CITATION.cff} (DOI to be assigned at the first tagged
release) live at the repository root.

\subsection*{Code availability}

The complete source code, test suite, examples, benchmark
reproducers, and Sphinx/Doxygen documentation source are at
\href{https://github.com/zhouzb79/exact_diagonalization_clean}%
{\nolinkurl{https://github.com/zhouzb79/exact_diagonalization_clean}}
under the MIT licence. The release commit corresponding to this
manuscript is the \texttt{phase-8-final} tag.

\section*{Acknowledgements}

We thank the developers of QuSpin~\cite{Weinberg2017QuSpin,
Weinberg2019QuSpin2}, $\mathcal{H}\Phi$~\cite{Wietek2018HPhi}, and
XDiag~\cite{Wietek2025XDiag} for setting the bar in this design
space; the MPI Forum and the NVIDIA cuBLAS / cuSPARSE / cuSOLVER /
NCCL teams for the building blocks; and the maintainers of the HDF
group's HDF5 library~\cite{Folk2011HDF5,HDF5website} for the
output-format substrate.

\paragraph{Funding information}
Funding sources to be filled in.

\paragraph{Author contributions}
Z.Z.\ designed and implemented the toolkit, wrote the test suite, and
wrote the manuscript.

\paragraph{Competing interests}
The author declares no competing interests.

% =========================================================================
\bibliography{refs}

\end{document}
